\section{مدل}

\begin{frame}{چه ساخته شد و مدل به چه چیزی نگاه می‌کند}
  \begin{itemize}
    \item ۱۰۰ ویژگی، ۱۳ خانواده‌ی مدل، غربال چهارمرحله‌ای، انتخاب نهایی یک مدل درختی
      کوانتایلی
    \item پروتکل: تقسیم زمانی با مبدأ غلتان؛ پنجره‌ی آزمون فقط یک‌بار باز شد
    \item مهم‌ترین ورودی‌ها: هویت سلف، برهم‌کنش روزهفته با شهر، حجم رزرو، نرخ دیروزِ همان
      سلول، شوک دیروز
  \end{itemize}
  \begin{figure}
    \centering
    \includegraphics[width=0.62\textwidth]{9.1_perm_importance_fa}
  \end{figure}
  \note{کلمه‌ی SHAP را نگو، بگو «سهم هر ویژگی در تصمیم مدل». نکته‌ی مهم اینکه این فهرست
  دقیقاً با یافته‌های بخش ۳ می‌خواند، یعنی مدل چیز عجیبی یاد نگرفته. شاهد:
  \lr{reports/phase9/9.2\_shap\_mean\_abs.csv}، \lr{reports/phase7/model\_comparison.md}.}
\end{frame}

\begin{frame}{۱۳ خانواده‌ی مدل امتحان شدند}
  \begin{itemize}
    \item از رگرسیون خطی ساده تا شبکه‌ی عصبی و مدل‌های بیزی — طیف کامل، نه فقط چند حدس
    \item هر خانواده حداقل یک نماینده تا انتهای مسیر آموزش و ارزیابی رفت
    \item خانواده‌ی درختی/بوستینگ (خ۲) در نهایت مدل برتر را داد؛ ترکیب و کالیبراسیون (خ۱۲،
      خ۱۳) لایه‌ای روی همه‌اند، نه یک رقیب مستقل
  \end{itemize}
  \begin{figure}
    \centering
    \includegraphics[width=0.92\textwidth]{slides_14b_model_families}
  \end{figure}
  \note{اگر استاد پرسید «چرا این‌همه مدل»: چون بدون مقایسه‌ی گسترده نمی‌شد ادعا کرد ساده‌ترین
  یا پیچیده‌ترین مدل کافی است — نتیجه‌ی بخش بعد (نتایج منفی) دقیقاً همین را نشان می‌دهد. شاهد:
  \lr{doc/WBS-phase7-modeling.md} بندهای 7.10--7.22.}
\end{frame}

\begin{frame}{نتایج منفی}
  \begin{itemize}
    \item از ۲۵ مدل فقط ۵ تا برد آماری معنادار داشتند
    \item شبکه‌های عصبی و مدل سطح فرد کمکی نکردند
    \item مدل برتر در اعتبارسنجی با مدل برتر در پنجره‌ی آزمون یکی نبود، و طبق پروتکل تغییر
      نکرد
  \end{itemize}
  \begin{figure}
    \centering
    \includegraphics[width=0.62\textwidth]{slides_15_funnel}
  \end{figure}
  \note{این اسلاید را حذف نکن. استادها روی همین قضاوت می‌کنند. شاهد: یافته‌های ۲۸، ۳۱، ۳۶؛
  ردیف ۴۹ decision\_log.}
\end{frame}

\begin{frame}{۵ مدلی که واقعاً بردند}
  \begin{itemize}
    \item از ۲۵ مدل، \lr{lightgbm\_quantile} (خ۲ درختی/بوستینگ) قهرمان نهایی است؛ نزدیک‌ترین
      رقیب یک مدل بیزی سلسله‌مراتبی (خ۸) است
    \item هر ۵ مدل برنده به‌طور معنادار از خط پایه \lr{B3} بهترند، با فاصله‌ای بین
      $0.0006$ تا $0.0012$ واحد \lr{pinball}
    \item رتبه‌بندی بر پایه‌ی \lr{pinball} نرخ است — تنها ستونی که بین هر دو گروه \lr{CPU}/\lr{GPU}
      این جدول مستقیماً قابل‌قیاس است
  \end{itemize}
  \begin{figure}
    \centering
    \includegraphics[width=0.86\textwidth]{slides_23_top_models_table}
  \end{figure}
  \note{شاهد: \lr{reports/phase7/model\_comparison.md}، ردیف‌های ۴۰ و ۴۶ decision\_log
  (انتخاب قهرمان و گسترش با ۴ قهرمان \lr{GPU}). \lr{bhm\_varying\_dispersion} چک‌لیست
  همگرایی خودش را کامل پاس نکرده (ردیف ۴۶) — اگر استاد پرسید چرا نزدیک‌ترین رقیب قهرمان
  نشد، همین را بگو.}
\end{frame}
